\section{A guitar in a sandstorm: forty years since Walsh}\label{S:wave-history}

Imagine leaving a guitar outside during a sandstorm. Each grain gives a string
an unpredictable impulse. The string carries the disturbance away, while new
grains strike elsewhere. How can we describe the motion when the individual
impacts are too numerous to follow? Walsh used this picture to introduce a
stochastic wave equation in his Saint-Flour lectures, delivered in 1984 and
published in 1986~\cite{walsh:86:introduction}*{p.~282, (1.14)}. Forty years
after their publication, the picture still provides a useful way into the
subject. We take these lectures as a starting point for this account. Earlier
work on stochastic wave equations lies outside its scope.

\begin{figure}
  \centering
  \begin{tikzpicture}
    \begin{scope}
      \clip[rounded corners=9pt]
      (-0.24\textwidth,-0.24\textwidth) rectangle
      (0.24\textwidth,0.24\textwidth);
      \node[inner sep=0pt,outer sep=0pt] (art) {
      \includegraphics[width=0.48\textwidth]{media/guitar-sandstorm/illustration.png}};
      % Soft white rings blend the image into the paper without changing the art.
      \begin{scope}
        \clip (art.south west) rectangle (art.north east);
        \foreach \k in {24,23,...,1} {
        \draw[white,opacity=0.12,line width=\k*0.65pt,
        rounded corners=9pt]
        (art.south west) rectangle (art.north east);
        }
      \end{scope}
      % The outer white frame also removes the square corner remnants.
      \draw[white,line width=3pt,rounded corners=9pt]
      (art.south west) rectangle (art.north east);
    \end{scope}
  \end{tikzpicture}
  \caption[A guitar in a sandstorm]{AI-generated illustration inspired by
    Walsh's guitar-in-a-sandstorm example. Artwork from Le Chen's
    \href{https://github.com/SPDEngine/SPDE-Visuals}{SPDE-Visuals collection},
    shared under the
    \href{https://creativecommons.org/publicdomain/zero/1.0/}{CC0 public-domain dedication}.
    This is an artistic illustration, not a
    numerical simulation. Walsh's worked example uses an infinite string to
    avoid boundary conditions. Companion videos and source information are
    included in \path{media/guitar-sandstorm/}.}
  \label{Fig:walsh-guitar}
\end{figure}


\paragraph{1986: from random impacts to a mathematical equation.}
A physical guitar has finite strings and endpoints, but Walsh makes a deliberate
simplification in the worked example: the string is infinite and initially at
rest, so boundary conditions do not obscure the random forcing. Centering the
impacts and replacing their fluctuations by space-time white noise gives an
additively forced wave equation. The differential notation is suggestive, but
the solution need not possess the derivatives appearing in it. The integral
formulation gives the equation its meaning. We must first specify what we mean
by a solution and the space in which it is defined. The guitar is a motivation
for this model, not a derivation of every stochastic wave equation considered
below.

\paragraph{1999: how far does the integration theory reach?}
The one-dimensional string is only the beginning. In higher dimensions the wave
propagator may itself be a measure, and spatially correlated noise calls for a
wider integration framework. Dalang's extension of the martingale-measure
integral~\cite{dalang:99:extending} made it possible to treat SPDEs with more
singular Green kernels under spectral integrability conditions. The question now
includes the interaction of propagation, spatial dimension and noise covariance.
These conditions determine whether the equation admits a random-field solution.

\paragraph{The 2000s and 2010s: from existence to growth.}
Once a solution is available, we can ask how its fluctuations develop over long
times. Multiplicative noise changes the problem: the forcing depends on the
current solution, so the disturbance and its random environment interact.
Moments measure the size of the fluctuations and the contribution of rare, large
values. Dalang--Mueller--Tribe~\cite{dalang.mueller.ea:08:feynman-kac-type}
developed probabilistic moment representations using Poisson processes and
random flights. Dalang--Mueller~\cite{dalang.mueller:09:intermittency} obtained
intermittency estimates for a three-dimensional model under their spatial
covariance assumptions. These results illustrate how the wave equation has its
own probabilistic methods, while the form of the noise continues to matter.

For one-dimensional space-time white noise,
Chen--Dalang~\cite{chen.dalang:15:moment} developed moment bounds and explicit
second-moment formulas, including rough initial data. Temporally correlated
forcing brings another change: random impacts at distinct times are no longer
independent. Balan--Conus~\cite{balan.conus:16:intermittency} studied the
resulting Skorohod moment growth under their covariance assumptions. For
white-in-time noise, Balan--Song~\cite{balan.song:19:second} related hyperbolic
and parabolic second-moment Laplace transforms and obtained an actual second
Lyapunov limit in their Riesz case. Construction, bounds of the correct growth
order and exact limits are successive questions, but a theorem answering one
need not answer the others.

\paragraph{The 2020s: exact rates and the meaning of multiplication.}
The same formal multiplicative equation can describe different stochastic
models. Skorohod integration uses Wick products; ordinary Stratonovich
approximation retains the self-contractions generated by multiplication of
smooth Gaussian fields. This distinction is central when the noise is static or
correlated in time. For time-independent Skorohod noise, Balan--Le Chen--Xia
Chen~\cite{balan.chen.ea:22:exact} obtained exact moment asymptotics.
Chen--Eisenberg~\cite{chen.eisenberg:23:interpolating} treated a related
fractional differential family, with time order below two and static noise.
Chen--Guo--Song~\cite{chen.guo.ea:24:moments} obtained exact second moments for
a white-noise fractional SPDE family and matching higher-moment growth orders in
its wave specialization. The latter are growth orders, not sharp higher-moment
constants.

On the Stratonovich side, Balan~\cite{balan:22:stratonovich} constructed
static-noise solutions and compared chaos expansions, while Xia
Chen--Hu~\cite{chen.hu:25:hyperbolic} established static homogeneous moment
asymptotics. For fractional temporal correlation, Xia Chen's author
version~\cite{chen:author2026:time-dependency} gives two-sided first-moment
bounds and a coarse second-moment upper bound. The source editions and all
model, data and normalization restrictions are detailed in
Section~\ref{S:literature-context}.

\paragraph{The present question.}
These works lead to three questions. How do we construct a solution? At what
rate do its moments grow? Can we identify the leading constant? Our problem is
the second moment of a multiplicative Stratonovich wave equation with fractional
temporal correlation and homogeneous spatial covariance. We give an upper bound
at the correct scale and a variational expression for the second-moment limit.
We do not obtain a numerical closed form or a sharp formula for every moment
order. The scientific review status is recorded separately. The following
sections specify the model and explain why its self-pairings and correlated
forcing require more than a transfer of the Skorohod results.
